% Burgers equation with viscous/artificial-viscosity shock
% Compiled from the discussion

\documentclass{mooiman_memo}

\memoSubject{Burgers Equation with Viscous and Artificial-Viscosity Shock}
\memoFrom{Jan Mooiman}
\memoDate{\today~\currenttime}
\memoTo{}

\begin{document}

\mooimantitle

\subsubsection*{Burgers Equation with Viscous and Artificial-Viscosity Shock}

That changes the interpretation substantially. With a delta formulation + Newton iteration, we can write down the discrete nonlinear problem much more precisely.

Suppose at time level n+1 you solve
\begin{align}
\mat{R}(\vec{u}^{n+1})=0
\end{align}
by Newton iteration.
At Newton iteration k,
\begin{align}
\mat{J}(\vec{u}(k))\delta \vec{u}=-\mat{R}(\vec{u}(k)),
\end{align}
followed by
\begin{align}
\vec{u}^{(k+1)}=\vec{u}^{(k)}+\delta \vec{u}.
\end{align}
Your matrix A is therefore the Newton Jacobian, not necessarily the physical diffusion operator itself. This is the important distinction that was missing in the previous derivation.

\subparagraph*{Your artificial-viscosity Jacobian}

You have
\begin{align}
A_{i,i - 1} & = \frac{1}{8} - c_{\Psi},
\\
A_{i,i} & = \frac{6}{8} + 2 c_{\Psi},
\\
A_{i,i+1} & = \frac{1}{8} - c_{\Psi}.
\end{align}
%
The artificial part of the Jacobian is
\begin{align}
c_{\Psi} [-1, 2, -1].
\end{align}
For the \textbf{delta variable},
\begin{align}
\delta u_{i-1} -2 \delta u_{i} + \delta u_{i+1}
=\Dx^2 \pdiff{\delta u}{x} + \frac{1}{12} \Dx^4 \pdiff[4]{\delta u}{x}.
\end{align}
Therefore the artificial-viscosity contribution to the Newton system is
\begin{align}
    \boxed{
    -c_{\Psi} \Dx^2 \pdiff{\delta u}{x} - \frac{c_{\Psi}}{12} \Dx^4 \pdiff[4]{\delta u}{x} + \ldots.
}
\end{align}
%
and this is a contribution to the \textbf{Jacobian}, not directly $\nu_{art}u_{xx}$.

\subparagraph*{The crucial point about your $c_E$}

You have
\begin{align}
    c_E = \frac{\pi}{2}c_\Psi^2
\end{align}
and
\begin{align}
    rhs_i = \frac{\pi}{4} c_\Psi^2 \Dx^2 \abs{u_{xx}}
\end{align}
%
Because this appears in the \textbf{Newton residual/RHS}, it represents a nonlinear residual contribution.
%
If the residual contains
\begin{align}
    \mat{R}_\Psi(\vec{u}) = \frac{\pi}{4} c_\Psi^2 \Dx^2 \abs{u_{xx}}
\end{align}

then its Newton linearisation is
\begin{align}
    \delta \mat{R}_\Psi = \frac{\pi}{4} c_\Psi^2 \Dx^2 {\rm sign}(u_{xx})\delta u_{xx}
\end{align}


Thus the Jacobian contribution associated with this term should be proportional to
\begin{align}
    \boxed{
    \frac{\pi}{4} c_\Psi^2 \Dx^2 {\rm sign}(u_{xx})\delta u_{xx}
    }
\end{align}

%
This is quite different from a conventional constant viscosity.

\subparagraph*{What your scheme is effectively doing}

Ignoring higher-order terms, your Newton equation contains artificial terms of the form
\begin{align}
    -c_{\Psi} \Dx^2 \delta u_{xx}
    + \frac{\pi}{4} c_\Psi^2 \Dx^2 {\rm sign}(u_{xx})\delta u_{xx}
\end{align}
%
Therefore the effective coefficient associated with the second derivative is approximately
\begin{align}
    \boxed{
    \nu_{\it eff} \propto    c_{\Psi} \Dx^2
    - \frac{\pi}{4} c_\Psi^2 \Dx^2 {\rm sign}(u_{xx})
    }
\end{align}
%
\textbf{up to the sign convention of your residual.}

That sign convention is important: depending on whether you define
\begin{align}
    R = F(u) - rhs
\end{align}
or
\begin{align}
    R = rhs - F(u)
\end{align}
the signs reverse.

So I would \textbf{not yet use this formula to choose $c_psi$.}

\subparagraph*{What I now need}

Since you have told me that the method is a delta/Newton formulation, the missing piece is very small. Please give me the actual Newton equation, something like
\begin{align}
\mat{A}  \delta = rhs;
\end{align}
and preferably the definition of $rhs$ immediately before this.
%
Then I can derive exactly:
\begin{align}
    &\boxed{
        R(u)
    }
    \qquad
    \boxed{
    J(u) = \pdiff{R(u)}{u}
    }
\end{align}
%
and consequently the continuous modified equation
\begin{align}
    u_u + uu_x = \nu_{\it eff}(u, u_x, u_{xx}) u_{xx} + HOT
\end{align}
%
That will also let us determine whether your current $c_E$ term is\textbf{ actually adding viscosity or inadvertently reducing it}, which is particularly relevant to the wiggles you observed for $a=1$, $b=-1$.
\end{document}
